\documentclass[tikz]{standalone}

\usepackage{pgfplots}
\pgfplotsset{compat=1.18}
\usetikzlibrary{arrows.meta, decorations.markings}

% Each edge pair keeps its color through every stage, so it stays visible which edge of
% the square becomes which circle on the torus.
\definecolor{rimcolor}{HTML}{0B5FA5}  % vertical edges -> cylinder rims -> short circle
\definecolor{seamcolor}{HTML}{C2570A} % horizontal edges -> the seam -> long circle
\definecolor{annotation}{HTML}{4A5560}
\definecolor{hairline}{HTML}{78828C}
\definecolor{meshgray}{HTML}{8B97A6}
\definecolor{tubeinside}{HTML}{8E959D}

% Camera shared by both 3D panels. The azimuth matters: viewed square-on, a tube's rims
% collapse to straight lines.
\def\viewh{-27}
\def\viewv{20}
\def\tuber{0.62}
\def\cylhalf{1.35}
\def\ringr{1.32}
\def\torusr{0.5}

% Lambert term for a tube, from its surface normal and a fixed light, fed to point meta
% raw. Arithmetic on it is not worth attempting: a comma inside point meta ends the key,
% and a comparison silently evaluates to a constant, either of which flattens the surface
% to a single tone. The colormap below does the clamping instead.
\def\cyldot{-0.66*cos(y) + 0.62*sin(y)}
\def\torusdot{-0.42*cos(y)*cos(x) - 0.66*cos(y)*sin(x) + 0.62*sin(y)}

\pgfplotsset{
  % quad edges double as the surface mesh, carrying the square's grid onto the rolled
  % up shapes. Leave point meta to auto-scale: pinning it flattens the shading.
  tube surface/.style={surf, shader=faceted, z buffer=sort, faceted color=meshgray},
  % colormap belongs on the axis. Set per plot it is ignored and the surface comes out
  % one flat tone, mesh and all.
  panel/.style={
    hide axis, axis equal image, anchor=center, view={\viewh}{\viewv},
    clip=false, scale only axis,
    % holding the first half flat clamps the unlit side, so only faces turned toward the
    % light brighten, the way a Lambert term with a floor would
    colormap={shade}{rgb(0cm)=(.640,.665,.695); rgb(1cm)=(.640,.665,.695);
      rgb(2cm)=(.885,.910,.945)},
  },
}

\tikzset{
  % one or two chevrons partway along a path: the usual shorthand for which edge is glued
  % to which, and in which direction
  midarrow/.style={postaction={decorate}, decoration={markings,
        mark=at position #1 with {\arrow{Stealth[length=8pt,width=7pt]}}}},
  midarrow/.default=.5,
  doublearrow/.style={postaction={decorate}, decoration={markings,
        mark=at position #1 with {\arrow{Stealth[length=8pt,width=7pt]}},
        mark=at position #1+.08 with {\arrow{Stealth[length=8pt,width=7pt]}}}},
  doublearrow/.default=.47,
  edge/.style={line width=2.2pt},
  cycle line/.style={line width=1.8pt},
}

% A labeled arrow marking one gluing step.
\newcommand{\steparrow}[3]{%
  \draw[annotation, -{Stealth[length=5pt]}, line width=.9pt] (#1,0) -- (#1+1.45,0);
  \node[#2, align=center, font=\small, anchor=south] at (#1+.725,.2) {#3};
}
\newcommand{\panelcaption}[2]{%
  \node[annotation, font=\small, anchor=north] at (#1,-2.35) {#2};
}

\begin{document}
\begin{tikzpicture}

  % ===== 1. the periodic plane =====
  % every vertical lattice line is a rim edge and every horizontal one a seam edge, so
  % tinting the whole lattice shows the identification repeating across the plane
  \foreach \i in {-1,0,1} {
    \draw[rimcolor!22, line width=.9pt] (1.3*\i,-1.95) -- (1.3*\i,1.95);
    \draw[seamcolor!22, line width=.9pt] (-1.95,1.3*\i) -- (1.95,1.3*\i);
  }
  \draw[annotation, -{Stealth[length=5pt]}, line width=.7pt] (-2.4,0) -- (2.4,0);
  \draw[annotation, -{Stealth[length=5pt]}, line width=.7pt] (0,-2.4) -- (0,2.4);
  % shifting this cell by 2pi in either direction lands on a copy of itself
  \fill[seamcolor!10] (0,0) rectangle (1.3,1.3);
  \draw[seamcolor, line width=1.5pt] (0,0) -- (1.3,0) (0,1.3) -- (1.3,1.3);
  \draw[rimcolor, line width=1.5pt] (0,0) -- (0,1.3) (1.3,0) -- (1.3,1.3);
  \node[font=\small, anchor=north] at (.65,-.22) {$2\pi$};
  \node[font=\small, anchor=west] at (1.46,.65) {$2\pi$};
  \panelcaption{0}{plane with $2\pi$ periodicity}

  % ===== 2. the fundamental domain and its two identifications =====
  \steparrow{2.55}{annotation}{one cell}

  \begin{scope}[xshift=4.75cm, yshift=-1.075cm]
    \draw[hairline!42, line width=.35pt, step=.5375] (0,0) grid (2.15,2.15);
    \draw[seamcolor, edge, midarrow] (0,0) -- (2.15,0);
    \draw[seamcolor, edge, midarrow] (0,2.15) -- (2.15,2.15);
    \draw[rimcolor, edge, doublearrow] (0,0) -- (0,2.15);
    \draw[rimcolor, edge, doublearrow] (2.15,0) -- (2.15,2.15);
    \node[rimcolor, font=\small, anchor=east] at (-.24,1.075) {$a$};
    \node[seamcolor, font=\small, anchor=north] at (1.075,-.16) {$b$};
  \end{scope}
  \panelcaption{5.825}{fundamental domain}

  % ===== 3. glue the horizontal pair -> cylinder =====
  \steparrow{7.6}{seamcolor}{glue\\top \& bottom}

  \begin{axis}[panel, at={(11.3cm,0cm)}, width=4.3cm, height=4.3cm]
    \addplot3[tube surface, samples=20, samples y=18,
      domain=-\cylhalf:\cylhalf, y domain=0:360, point meta={\cyldot}]
      ({x}, {\tuber*cos(y)}, {\tuber*sin(y)});
    % looking into the open near end you see the far inner wall. pgfplots shades back
    % faces as if they were outside, so fill that disc with the shade the inside carries.
    \addplot3[fill=tubeinside, draw=none, samples=40, domain=0:360]
      ({-\cylhalf}, {\tuber*cos(x)}, {\tuber*sin(x)});
    % the glued pair is now a single seam running the length of the tube
    % swept from the far end back, so the chevron points the same way as the square's
    \addplot3[seamcolor, cycle line, midarrow=.45, samples=2, domain=\cylhalf:-\cylhalf]
      ({x}, {0}, {\tuber});
    % the rims are still open: they are the two vertical edges of the square. The whole
    % near rim bounds the opening we look into, but only 68:248 of the far one clears the
    % tube body; past that its normal turns away from the camera.
    \addplot3[rimcolor, cycle line, doublearrow=.52, samples=40, domain=0:360]
      ({-\cylhalf}, {\tuber*cos(x)}, {\tuber*sin(x)});
    \addplot3[rimcolor, cycle line, doublearrow=.30, samples=40, domain=68:248]
      ({\cylhalf}, {\tuber*cos(x)}, {\tuber*sin(x)});
  \end{axis}
  \panelcaption{11.3}{cylinder}

  % ===== 4. bend it round and glue the rims -> torus =====
  \steparrow{13.5}{rimcolor}{glue\\rim to rim}

  \begin{axis}[panel, at={(17.1cm,0cm)}, width=4.7cm, height=4.7cm]
    \addplot3[tube surface, samples=40, samples y=16,
      domain=0:360, y domain=0:360, point meta={\torusdot}]
      ({(\ringr+\torusr*cos(y))*cos(x)}, {(\ringr+\torusr*cos(y))*sin(x)}, {\torusr*sin(y)});
    % the seam closed into the long way round; the two rims fused into one short circle
    \addplot3[seamcolor, cycle line, midarrow=.62, samples=60, domain=0:360]
      ({\ringr*cos(x)}, {\ringr*sin(x)}, {\torusr});
    % likewise 24:203 for a meridian sitting at azimuth -18, on the right of the ring
    \addplot3[rimcolor, cycle line, doublearrow=.62, samples=40, domain=24:203]
      ({(\ringr+\torusr*cos(x))*cos(-18)}, {(\ringr+\torusr*cos(x))*sin(-18)},
       {\torusr*sin(x)});
    \node[seamcolor, font=\small, anchor=south] at (axis cs:0,0,1.05) {$b$};
    \node[rimcolor, font=\small, anchor=north west] at (axis cs:2.0,-.65,-.35) {$a$};
  \end{axis}
  \panelcaption{17.1}{torus}

\end{tikzpicture}
\end{document}
